\section{A uniform small-solution estimate for the lift}

\begin{lemma}[Uniform small-solution interface]\label{lem:savin-family}
Fix $N$, $\lambda_0>0$, $\Lambda_0<\infty$, $\rho>0$ and $K<\infty$.
Consider autonomous continuous degenerate-elliptic operators $F(X,q)$
with $F(0,q)=0$. Suppose they are $C^2$ on the common jet ball
$|X|,|q|<\rho$, with common $C^2$ bound $K$, and their matrix
derivatives have eigenvalues in $[\lambda_0,\Lambda_0]$ there.
There are common $\sigma>0$, $\alpha>0$ and $C<\infty$ such that
any continuous viscosity solution of $F(D^2Z,DZ)=0$ on $B_1$ with
$\|Z\|_{L^\infty}\leq\sigma$ is $C^{2,\alpha}$ on $B_{1/2}$ with
norm at most $C$. The threshold can be reduced to obtain any fixed
sufficiently small bound on the first and second derivatives.
\end{lemma}
\begin{proof}
This is the common-data form of \cite[Theorem 1.3 and
Proposition 4.1, pp. 571--573]{Savin}. The proposition's contradiction
argument explicitly allows varying $F_k$ under a common second-derivative
bound, extracts $F_k,DF_k$, and obtains a uniformly elliptic
constant-coefficient limiting equation. Its quadratic approximation
and iteration constants depend only on the common local data.
The lower-order and space variables in that source are absent here.
Its H2 uses the matrix operator norm; the eigenvalue assumptions above
give its constants $\lambda_0$ and $N\Lambda_0$. Reducing the common
jet radius gives the last assertion. No gradient bound for the unknown
solution is assumed.
\end{proof}

\begin{lemma}[Coarse-cylinder estimate]\label{lem:coarse-cylinder}
Fix $n,L,b_0>0$. There are $d_*,k_*>0$ and $C<\infty$, depending
only on these data, with the following property. Suppose an actual
affine lift $U$ of a weak $p$-harmonic solution, $p\geq3$,
is defined on a neighborhood of an $\eta$-fitted cylinder
$Z_\eta(z,H;P)$, with slope $P=ce_N$ after rotation, and
\begin{equation}\label{eq:coarse-hyp}
 \begin{gathered}
 c\in[1/2,L+1],\quad A=1-c^{-2}>0,\quad
 \epsilon=(p-2)^{-1}\leq1,\quad e=\epsilon A\leq1,\\
 \xi=(a,b),\quad |a|^2+b^2=1,\quad b\geq b_0,\qquad 0<\eta\leq d_*e.
 \end{gathered}
\end{equation}
Then $P_{\mathrm{new}}=DU(z)$ exists and is nonzero, and, whenever
$0<\lambda\leq e$, $0<r\leq k_*\ell$ with
$\ell=(H/4)\sqrt{\eta/e}$, one has
\begin{equation}\label{eq:cylinder-estimate}
 \sup_{h:\,z+h\in Z_\lambda(z,r;P_{\mathrm{new}})}
 |U(z+h)-U(z)-P_{\mathrm{new}}\cdot h|\leq Ce r^2/H.
\end{equation}
\end{lemma}
\begin{proof}
Put $T_e=\operatorname{diag}(\sqrt e I,1)$ and set
\[
 U(z+\ell T_e(x,t))=U(z)+c\ell[t+eW(x,t)].
\]
For $d_*\leq1$ this ellipsoid lies strictly in the fitted cylinder:
its transverse radius is $H\sqrt\eta/4$ and its axial radius is at
most $H/4$. The value fit gives
\begin{equation}\label{eq:W-small}
 \|W\|_{L^\infty(B_1)}\leq(4/c)\sqrt{\eta/e}.
\end{equation}
At constrained tests its equation is
\[
 D^2W[(D_xW,1+eW_t)]+B(DW)\tr(S_e\Pi S_eD^2W)=0,
\]
where $S_e=\operatorname{diag}(I,\sqrt e)$ and
\[
 B(q)=1+\epsilon(2q_t+|q_x|^2)+\epsilon^2Aq_t^2.
\]
At those tests $B=|\Pi(\sqrt e q_x,1+e q_t)|^2/A\geq0$.
Again replace $B$ globally by its nonnegative part. To obtain a
proportional, rather than absolute, derivative bound include an
amplitude parameter $d\in[0,1]$:
\begin{align*}
 F_d(X,q)&=X[(d q_x,1+edq_t)]
           +\max\{B_d(q),0\}\tr(S_e\Pi S_eX),\\
 B_d(q)&=1+\epsilon(2dq_t+d^2|q_x|^2)
                          +\epsilon^2Ad^2q_t^2.
\end{align*}
Allowing $\epsilon,A,e,d$ independently in $[0,1]$ only enlarges
the family. All members are continuous and degenerate elliptic on
all jets. On $|q|\leq1/8$ one has $B_d\geq3/4$, so the max switch
is absent. They vanish at $X=0$ and have common local $C^2$ bounds.
At $q=0$ their matrix form is
\[
 Q(h,s)=|h|^2+(1+e)s^2-(a\cdot h+b\sqrt e s)^2,
 \qquad \frac{b_0^2}{3}(|h|^2+s^2)\leq Q(h,s)
                    \leq2(|h|^2+s^2).
\]
Here is a dimension-free certificate for this actual vector form.
Set $t=\sqrt e$ and
$g=|a|^2|h|^2-\langle a,h\rangle^2\geq0$, by Cauchy--Schwarz.
The unit relation $|a|^2+b^2=1$ gives
\[
 Q=|bh-tsa|^2+s^2+g
\]
and
\[
 \begin{split}
 3Q-b^2(|h|^2+s^2)
 ={}&|bh-2tsa|^2+|bh-tsa|^2+3g\\
    &+(3-b^2-2t^2|a|^2)s^2\geq0.
 \end{split}
\]
Indeed $3-b^2-2t^2|a|^2\geq1+b^2$ for $t^2\leq1$.
The upper bound follows by dropping the negative square in $Q$.
This works on any real inner product space and includes $a=0$.
The supporting Lean theorem \texttt{physical\_uniform\_bounds}
certifies precisely these bounds with $t=\sqrt e$, the unit relation,
$0\leq e\leq1$ and $b_0^2\leq b^2$.

We quantify the common neighborhood as well. Write
\[
 v(q)=(dq_x,1+edq_t),\qquad
 M_d(q)=v(q)v(q)^T+B_d(q)S_e\Pi S_e.
\]
For $|q|\leq r\leq1/8$, one has $|v(q)-e_N|\leq r$,
$|B_d(q)-1|\leq2r+2r^2$, and $\|S_e\Pi S_e\|\leq1$. Hence
\[
 \|M_d(q)-M_d(0)\|\leq4r+3r^2<5r.
\]
Thus the explicit common radius
\[
 \rho=\min\{1/8,b_0^2/30\}
\]
gives $B_d\geq3/4$ and eigenvalues in $[b_0^2/6,3]$ throughout
$|q|\leq\rho$, including all parameter endpoints. For every positive
semidefinite matrix $Y$, the operator-norm convention gives
\[
 \frac{b_0^2}{6}\|Y\|\leq\tr(M_d(q)Y)
                  \leq3\tr Y\leq3N\|Y\|.
\]
Savin's ellipticity constants can therefore be taken as
$\lambda=b_0^2/6$ and $\Lambda=3N$. On $|X|,|q|\leq\rho$ the
max switch is absent, so the coefficients are polynomial in $q$;
$D_qM_d$ and $D_q^2M_d$ are uniformly bounded. Consequently
$D^2_{Xq}F_d=D_qM_d$ and
$D^2_{qq}F_d=\tr(D_q^2M_dX)$ have common bounds depending only on
$N,b_0$. Together with $F_d(0,q)=0$, these verify every hypothesis
of Lemma~\ref{lem:savin-family} with one common $K$. The square
root in $S_e$ causes no loss: parameter derivatives are not required.

Let $a_0=\|W\|\leq\sigma$, with the common threshold from
Lemma~\ref{lem:savin-family}. If $a_0=0$ the estimate is immediate.
Otherwise take $d=a_0/\sigma$ and $Z=W/d$. Dividing the equation by
$d$ gives exactly $F_d(D^2Z,DZ)=0$, and $\|Z\|=\sigma$. Thus
\[
 \|DW\|_{B_{1/2}}+\|D^2W\|_{B_{1/2}}\leq C a_0.
\]
Choose $d_*$ so that \eqref{eq:W-small} is below the threshold.
Scaling back yields
\[
 \|T_eD^2UT_e\|\leq\frac{ce}{\ell}C\frac4c\sqrt{\eta/e}
                  \leq C'e/H
\]
on the half ellipsoid. By decreasing the threshold its derivative
conclusion also gives $c(1+eW_t)\geq c/2$. Hence
\[
 P_{\mathrm{new}}=c(\sqrt e D_xW(0),1+eW_t(0))\neq0,
 \qquad |T_e^{-1}N_{P_{\mathrm{new}}}|\leq C.
\]
For $h$ in its $\lambda$ cylinder, $\lambda\leq e$ implies
$|T_e^{-1}h|\leq Cr+\sqrt{\lambda/e}r\leq C'r$.
Choose $k_*$ so the whole segment lies in the half ellipsoid.
Taylor's formula proves \eqref{eq:cylinder-estimate}.
\end{proof}

\section{Infinite effective viscosity and the quadratic comparison}

It remains to exclude $\nu_i\to+\infty$. Then $A_i>0$ eventually.
Define
\[
 R_i=\nu_i^2,\quad H_i=R_ir_i,\quad
 \kappa_i=64(4R_i)^{\gamma_i},\quad \eta_i=\kappa_i\lambda_i.
\]
Since $A_i,\epsilon_i\leq1$, $R_i\leq\lambda_i^{-2}$, and
\eqref{eq:scale-separation} gives $R_is_i\to0$. Centers within
$2H_i$ of $z_i$ retain at least half its boundary distance. At all
admissible radii up to $H_i$, \eqref{eq:running-max} gives
\[
 \beta\leq b(4R_is_i)\leq8\lambda_i(4R_i)^{\gamma_i}<\eta_i.
\]
Thus $\eta_i$ is good throughout these comparisons. Also
$\eta_i\leq C\lambda_i^{1-2\gamma_i}\to0$ and
$\kappa_i/\nu_i\to0$. A fixed multiple of $\lambda_i$ is not
asserted good on these growing boxes.

Let $Q_{i,0}$ be an $\eta_i$-fit at the base radius $r_i$, and
let $P_i$ be the original $\Lambda\lambda_i$-fit. Their common
inscribed ball gives
$|P_i-Q_{i,0}|\leq C\kappa_i\sqrt{\lambda_i}$, hence the same
order for their axis angle. An equal-parameter norm comparison
cannot be applied to this pair. Instead take a point on the base
axis at distance $t_ir_i$ where $t_i=c\kappa_i^{-1/2}$ and $c>0$
is sufficiently small and fixed. Its transverse displacement from
the coarse axis is at most
$Ct_i\kappa_i\sqrt{\lambda_i}r_i\leq\sqrt{\eta_i}r_i$;
the axial margins also hold. Evaluating both fits there gives
\[
 |(P_i-Q_{i,0})\cdot N_{P_i}|
       \leq C\kappa_i^{3/2}\lambda_i.
\]
Converting axial projection to norm costs the square of the angle,
at most $C\kappa_i^2\lambda_i$. Therefore
\begin{equation}\label{eq:unequal-norm}
 \bigl||P_i|-|Q_{i,0}|\bigr|\leq C\kappa_i^2\lambda_i.
\end{equation}
Now chain $\eta_i$-fits through dyadic radii from $r_i$ to $H_i$.
These comparisons have equal parameters and common slope bounds,
so the source's quadratic norm comparison is applicable. A final
coarse fit $Q_i$ satisfies
\begin{equation}\label{eq:coarse-norm}
 \bigl||Q_i|-c_i\bigr|\leq
 C\{\kappa_i^2\lambda_i+\eta_i(1+\log\nu_i)\}.
\end{equation}
Put $c_i'=|Q_i|$, $A_i'=1-(c_i')^{-2}$. The map
$c\mapsto1-c^{-2}$ is uniformly Lipschitz on $[1/2,L+1]$, so
\[
 \frac{|A_i'-A_i|}{A_i}
 \leq C\epsilon_i
   \frac{\kappa_i^2+\kappa_i(1+\log\nu_i)}{\nu_i}
 \longrightarrow0.
\]
Indeed $\kappa_i=64\,4^{\gamma_i}\nu_i^{2\gamma_i}$,
$\gamma_i\to0$ and $\nu_i\to\infty$; eventually
$4\gamma_i<1/2$. Thus $A_i'>0$ and
$e_i'=\epsilon_i A_i'\sim\nu_i\lambda_i$.
This proves the needed comparison before substituting the coarse
parameter into the small-solution lemma.

The coarse fit has
$\eta_i/e_i'\sim\kappa_i/\nu_i\to0$. Its added-coordinate
component obeys
$b_i'\geq(1-\sqrt{\eta_i})/(L+1)$ by
\eqref{eq:slope-bounds}. Its ellipsoid scale is
\[
 \frac{\ell_i}{r_i}
 =\frac{R_i}{4}\sqrt{\eta_i/e_i'}
 \sim\frac14\nu_i^{3/2}\sqrt{\kappa_i}\longrightarrow\infty.
\]
Also $\lambda_i\leq e_i'$ eventually. Estimate
\eqref{eq:cylinder-estimate} at $r_i$ therefore has error, divided
by $\lambda_ir_i$, bounded by
\[
 \frac{Ce_i'r_i}{\lambda_iH_i}\sim\frac C{\nu_i}\longrightarrow0.
\]
Its nonzero slope is the actual $DU_i(z_i)$. It makes
$\lambda_i$ good, contradicting its selection. Every domain used
lies in the original solution region by
\eqref{eq:scale-separation}. We have excluded all branches of the
bad sequence and proved the required power bound.
